\documentclass[12pt,a4paper]{article}

\usepackage[spanish,es-nodecimaldot]{babel}
\usepackage[utf8]{inputenc}
\usepackage[T1]{fontenc}
\usepackage{lmodern}
\usepackage[a4paper,top=2.5cm,bottom=2.5cm,left=2.7cm,right=2.7cm]{geometry}
\usepackage{amsmath}
\usepackage{graphicx}
\usepackage{booktabs}
\usepackage{array}
\usepackage{enumitem}
\usepackage{xurl}
\usepackage[colorlinks=true,linkcolor=blue,urlcolor=blue,citecolor=blue]{hyperref}

\setlist{itemsep=2pt,topsep=4pt}
\setlength{\parskip}{0.45em}
\setlength{\parindent}{0pt}

\title{Servicio seguro de almacenamiento de archivos multimedia\\
\large Memoria de implementación -- primera fase}
\author{Integrantes: \textit{completar}}
\date{Curso 2026--2027}

\begin{document}
\maketitle

\begin{abstract}
Este trabajo implementa una aplicación de escritorio para cifrar y recuperar archivos mediante un esquema híbrido AES y RSA. Cada archivo se cifra con una clave AES-128 aleatoria en modo GCM; la clave AES se cifra con RSA-OAEP y la clave privada RSA se protege con AES-GCM mediante una clave derivada de una contraseña. La aplicación gráfica permite almacenar y consultar el contenido cifrado y las claves protegidas en una base de datos SQLite local. Esta memoria describe la arquitectura implementada, las decisiones de almacenamiento, el uso de la aplicación, las comprobaciones realizadas y las funcionalidades que quedan pendientes respecto a la segunda fase del enunciado.
\end{abstract}

\tableofcontents
\newpage

\section{Introducción}

El objetivo de la primera fase es proteger archivos multimedia con una clave simétrica aleatoria distinta para cada archivo y evitar que las claves AES almacenadas queden expuestas. Para ello se combinan algoritmos simétricos y asimétricos:

\begin{itemize}
    \item AES-128-GCM protege el contenido de cada archivo y permite detectar alteraciones.
    \item RSA-OAEP protege la clave AES aleatoria asociada a ese archivo.
    \item AES-128-GCM, con una clave derivada de la contraseña, protege la clave privada RSA.
    \item SQLite conserva el contenido cifrado y el material criptográfico protegido de la aplicación gráfica.
\end{itemize}

La aplicación está implementada en Python y cuenta con una interfaz de escritorio en Tkinter. El diseño separa la interfaz, la criptografía, el almacenamiento y los servicios que conectan esas piezas.

\section{Alcance del trabajo}

La versión actual implementa un flujo local de cifrado y descifrado para archivos. Trabaja sobre bytes, por lo que puede cifrar imágenes, audio, vídeo y otros tipos de archivo; no se limita a una extensión concreta.

La aplicación gráfica permite crear cuentas locales y autenticar al usuario mediante la contraseña que desbloquea su clave privada RSA. Cada cuenta posee su propio par RSA y sus propios archivos. La compartición segura entre usuarios todavía no está implementada.

\section{Tecnologías y dependencias}

\begin{itemize}
    \item \textbf{Python 3.9 o posterior}: lenguaje de implementación.
    \item \textbf{PyCryptodome}: implementación de AES, RSA y generación de bytes aleatorios criptográficamente seguros. La dependencia del proyecto está declarada en \texttt{requirements.txt}.
    \item \textbf{Tkinter}: interfaz gráfica, mediante la biblioteca estándar de Python.
    \item \textbf{SQLite}: base de datos local, mediante el módulo \texttt{sqlite3} de la biblioteca estándar; no requiere un servidor independiente.
    \item \textbf{unittest}: pruebas automatizadas de la capa de almacenamiento y del flujo criptográfico.
\end{itemize}

La interfaz gráfica también requiere que Tkinter esté instalado en el intérprete de Python. Algunas distribuciones de Linux ofrecen Tkinter como paquete del sistema separado.

\section{Arquitectura}

El proyecto se organiza en módulos para evitar mezclar la interfaz y las consultas SQL con las operaciones criptográficas:

\begin{description}[style=nextline,leftmargin=0pt]
    \item[\texttt{crypto/aes.py}]
    Genera claves AES-128 y proporciona funciones para cifrar y descifrar bytes o archivos con AES-GCM.
    \item[\texttt{crypto/rsa.py}]
    Genera un par RSA de 2048 bits y cifra o descifra claves simétricas con RSA-OAEP.
    \item[\texttt{crypto/passwords.py}]
    Deriva claves desde contraseñas con PBKDF2-HMAC-SHA256 y protege o recupera la clave privada RSA.
    \item[\texttt{database/database.py}]
    Inicializa el esquema SQLite y proporciona operaciones de lectura y escritura parametrizadas.
    \item[\texttt{services/file\_service.py}]
    Coordina cifrado, RSA y persistencia en la base de datos, sin depender de los widgets de Tkinter.
    \item[\texttt{main.py}]
    Crea la interfaz gráfica, valida la selección de archivos y solicita al servicio las operaciones.
    \item[\texttt{encrypt.py} y \texttt{decrypt.py}]
    Conservan un flujo alternativo de línea de comandos que guarda el contenido cifrado como archivos externos.
\end{description}

\section{Diseño criptográfico}

\subsection{Cifrado de cada archivo}

Para cada archivo nuevo, la aplicación genera una clave aleatoria de 16 bytes:

\[
16\ \text{bytes} \times 8 = 128\ \text{bits}.
\]

Esa clave se utiliza una sola vez para cifrar el contenido con AES en modo Galois/Counter (GCM). GCM produce el texto cifrado y una etiqueta de autenticación. La implementación guarda juntos el nonce, la etiqueta y el texto cifrado. Al descifrar, la etiqueta se verifica antes de escribir el contenido recuperado; si la autenticación falla, la operación produce un error.

\subsection{Protección de la clave AES}

RSA se utiliza para cifrar la clave AES, no el archivo completo. La clave pública cifra los 16 bytes de la clave AES con RSA-OAEP. Durante la recuperación, la clave privada correspondiente descifra ese material. Esta combinación permite usar el cifrado simétrico para archivos y el asimétrico para transportar la clave de cada archivo.

\subsection{Protección de la clave privada RSA}

La contraseña de la cuenta se utiliza exclusivamente para proteger y desbloquear su clave privada RSA; no cifra directamente los archivos ni sus claves AES. La contraseña no se almacena. Para proteger la clave privada RSA:

\begin{enumerate}
    \item Se genera un salt aleatorio de 16 bytes.
    \item PBKDF2-HMAC-SHA256 deriva 16 bytes de clave, con 200\,000 iteraciones en la implementación actual.
    \item La clave derivada cifra la representación DER de la clave privada con AES-128-GCM.
    \item Se conserva el salt junto con el nonce, la etiqueta de autenticación y el texto cifrado de la clave privada.
\end{enumerate}

El salt no es secreto. Permite repetir la derivación durante el descifrado. Una contraseña equivocada deriva otra clave y la verificación de GCM falla. Al iniciar sesión, la clave privada se desbloquea una vez y se conserva en memoria durante esa sesión; la interfaz no vuelve a solicitar la contraseña para cada operación.

\section{Proceso de cifrado y descifrado}

\subsection{Cifrado}

Desde la interfaz gráfica, el proceso es el siguiente:

\begin{enumerate}
    \item El usuario inicia sesión con su cuenta y desbloquea su clave privada RSA una vez.
    \item Elige un archivo que se cifra con una clave AES aleatoria nueva y AES-GCM.
    \item La clave AES se cifra con RSA-OAEP y la clave pública de la cuenta activa.
    \item El texto cifrado y la clave AES cifrada se guardan en SQLite junto con metadatos del archivo.
\end{enumerate}

\subsection{Descifrado}

Para recuperar un registro desde la biblioteca SQLite:

\begin{enumerate}
    \item Se selecciona un registro de la cuenta autenticada. Su clave privada ya está desbloqueada en memoria.
    \item RSA descifra la clave AES asociada al registro.
    \item AES-GCM descifra y autentica el contenido.
    \item Solo tras superar la autenticación, se escribe el archivo en la ubicación elegida.
\end{enumerate}

La aplicación también mantiene una pestaña para descifrar archivos externos \texttt{.enc} y \texttt{.key.enc}. La clave AES debe haberse cifrado para la clave pública de la cuenta activa. Esos archivos no se importan automáticamente a la biblioteca SQLite.

\section{Base de datos SQLite}

La interfaz crea automáticamente la base de datos local \texttt{secure\_storage.sqlite3}. Durante el desarrollo se ubica en \texttt{data/}. Cuando se ejecuta como aplicación empaquetada, el código está preparado para usar la carpeta del usuario \texttt{\%LOCALAPPDATA\%\textbackslash CS2026}, para no intentar escribir junto al ejecutable.

\subsection{Esquema}

\begin{center}
\small
\begin{tabular}{p{0.21\linewidth}p{0.34\linewidth}p{0.37\linewidth}}
\toprule
\textbf{Tabla} & \textbf{Campos principales} & \textbf{Contenido y propósito}\\
\midrule
\texttt{users} &
\texttt{id}, \texttt{username}, \texttt{public\_key}, \texttt{encrypted\_private\_key}, \texttt{created\_at} &
Almacena una cuenta y su par RSA. La clave privada está cifrada con AES-GCM; la contraseña no se guarda.\\
\addlinespace
\texttt{encrypted\_files} &
\texttt{id}, \texttt{owner\_id}, \texttt{original\_name}, \texttt{encrypted\_content}, \texttt{encrypted\_aes\_key}, \texttt{created\_at}, \texttt{size\_bytes} &
Almacena el propietario, el contenido cifrado y la clave AES cifrada como BLOB, además de metadatos.\\
\bottomrule
\end{tabular}
\end{center}

Las relaciones se protegen con claves foráneas de SQLite. Las consultas que insertan datos usan parámetros SQL, y las escrituras se realizan en transacciones. La biblioteca filtra por el identificador del usuario autenticado.

El contenido original y la clave AES en claro no se guardan en la base de datos. Sin embargo, SQLite no cifra automáticamente el fichero completo: los nombres, tamaños y fechas se conservan como metadatos legibles, y el archivo de base de datos debe mantenerse en un lugar protegido. Los archivos se leen y cifran íntegramente en memoria, por lo que esta primera versión está orientada a archivos de tamaño moderado.

Al abrir una base de datos de la versión anterior, el antiguo par RSA y sus archivos se migran a una cuenta local llamada \texttt{usuario}. Para iniciar sesión en esa cuenta se necesita la contraseña que protegía el par anterior.

\section{Interfaz y manual de usuario}

\subsection{Preparación}

Desde la carpeta raíz del proyecto, se crea y activa un entorno virtual.

\textbf{Linux o macOS:}
\begin{verbatim}
python3 -m venv .venv
source .venv/bin/activate
\end{verbatim}

\textbf{Windows PowerShell:}
\begin{verbatim}
py -m venv .venv
.\.venv\Scripts\Activate.ps1
\end{verbatim}

Con el entorno activado, se instala PyCryptodome y se inicia la interfaz:

\begin{verbatim}
python -m pip install -r requirements.txt
python main.py
\end{verbatim}

\subsection{Pestaña «Cifrar»}

Al abrir la aplicación, se elige una cuenta y se introduce su contraseña. Si todavía no hay cuentas, se puede crear una; durante el alta se genera un par RSA propio y se protege la clave privada con la contraseña elegida.

\begin{enumerate}
    \item Seleccionar el archivo original con el botón de exploración.
    \item Pulsar \textbf{Cifrar y guardar en SQLite}.
    \item La aplicación cifra para la cuenta activa e informa del identificador del registro guardado.
\end{enumerate}

El contenido cifrado y las claves se guardan en la base de datos local. No se genera un archivo cifrado externo en este flujo gráfico.

\subsection{Pestaña «Biblioteca SQLite»}

La biblioteca muestra el identificador, nombre original, fecha de almacenamiento y tamaño de los registros de la cuenta activa. Se puede actualizar la lista y elegir un registro para descifrarlo. La clave privada desbloqueada al iniciar sesión recupera el archivo; solo hay que seleccionar el destino. La interfaz confirma antes de sobrescribir un destino que ya exista.

\subsection{Pestaña «Descifrar»}

Esta pestaña sirve para archivos externos. Se eligen el archivo cifrado, su clave AES cifrada \texttt{.key.enc} y la ruta de salida. Se utiliza la clave privada de la cuenta activa y no se vuelve a pedir la contraseña. El modo de línea de comandos es independiente y recibe \texttt{--password} en cada ejecución; solo sus archivos cifrados con la clave pública de la cuenta activa se pueden abrir desde esta pestaña.

\section{Uso desde la línea de comandos}

El modo de consola es compatible con el flujo de ficheros externos.

\textbf{Cifrar:}
\begin{verbatim}
python encrypt.py /ruta/al/archivo.jpg --password "mi-contraseña"
\end{verbatim}

Por defecto, los archivos cifrados y las claves se crean en \texttt{data/}. Se puede cambiar la ubicación de salida:

\begin{verbatim}
python encrypt.py /ruta/al/archivo.jpg \
  --password "mi-contraseña" \
  --output-dir /ruta/de/salida
\end{verbatim}

\textbf{Descifrar:}
\begin{verbatim}
python decrypt.py \
  --encrypted-file data/archivo.enc \
  --key-file data/archivo.key.enc \
  --password "mi-contraseña" \
  --output /ruta/recuperada/archivo.jpg
\end{verbatim}

Si la clave privada protegida no está junto al archivo cifrado ni en \texttt{data/}, se indica su ruta mediante la opción \texttt{--private-key}.

\section{Pruebas y comprobaciones}

La suite \texttt{tests/test\_file\_service.py} contiene pruebas automatizadas para:

\begin{itemize}
    \item Cifrar un archivo, persistir el contenido y recuperar bytes idénticos al original.
    \item Crear cuentas, desbloquear sus claves privadas y rechazar contraseñas incorrectas.
    \item Aislar los archivos de cada usuario y evitar que otra cuenta los descifre.
    \item Migrar el esquema anterior de un solo par RSA conservando los archivos cifrados.
    \item Rechazar nombres de usuario duplicados.
\end{itemize}

Además, se ha comprobado manualmente la inicialización de la ventana Tkinter con la pestaña de biblioteca SQLite y se ha validado que el flujo de línea de comandos de cifrado/descifrado sigue siendo compatible con ficheros externos.

\section{Seguridad, límites y trabajo futuro}

\subsection{Precauciones}

\begin{itemize}
    \item Deben conservarse la base de datos SQLite y la contraseña para recuperar archivos guardados desde la interfaz.
    \item Perder la contraseña o la clave privada protegida hace irrecuperables los archivos asociados.
    \item La base de datos contiene datos cifrados, pero no es una base de datos cifrada íntegramente. Se recomienda proteger el acceso al dispositivo y realizar copias de seguridad privadas.
    \item El nombre original, el tamaño y la fecha de almacenamiento no se cifran.
    \item \texttt{.gitignore} excluye datos locales y claves, pero no elimina archivos que Git ya siga ni información publicada anteriormente. Antes de compartir el repositorio debe revisarse su historial y su estado.
\end{itemize}

\subsection{Estado respecto al enunciado}

La primera fase está implementada a nivel local: cifrado AES-128 por archivo, envoltura RSA de la clave AES, protección de la clave privada, cuentas locales con inicio de sesión y persistencia SQLite en la interfaz. La aplicación también ofrece un flujo por consola independiente.

Queda pendiente la segunda fase del enunciado:

\begin{enumerate}
    \item Compartición segura, cifrando la clave AES del archivo con la clave pública del destinatario y registrando el permiso correspondiente.
\end{enumerate}

También queda como tarea de entrega empaquetar y probar la aplicación en Windows. La ruta de datos del modo congelado está preparada para \texttt{\%LOCALAPPDATA\%}, pero esta memoria no afirma que ya se haya generado o validado un ejecutable.

\end{document}
